\documentclass[11pt,a4paper]{article}
\usepackage[utf8]{inputenc}
\usepackage[T1]{fontenc}
\usepackage{geometry}
\geometry{margin=1in}
\usepackage{hyperref}
\usepackage{enumitem}
\usepackage{xcolor}
\usepackage{titlesec}
\usepackage{listings}

% Design Colors
\definecolor{rwanda-blue}{HTML}{00A1DE}
\definecolor{dark-gray}{HTML}{333333}

\hypersetup{
    colorlinks=true,
    linkcolor=rwanda-blue,
    urlcolor=rwanda-blue,
    pdftitle={Software Engineer Internship Submission - Umuganda System},
}

\titleformat{\section}{\Large\bfseries\color{dark-gray}}{}{0em}{}[\titlerule]
\titleformat{\subsection}{\large\bfseries\color{dark-gray}}{}{0em}{}

\title{Software Engineer Internship Submission Strategy\\ \large Technical Project Case Study: Umuganda Management System}
\author{Andre Mugabo}
\date{\today}

\begin{document}

\maketitle

\section{Introduction}
This document serves as a technical portfolio and submission strategy for the Software Engineer Intern position at Fortis Green. It outlines the select project, technical architecture, problem-solving lifecycle, and a structured presentation script for technical evaluation.

\section{Project Selection: Umuganda Management System}
The \textbf{Umuganda Management System} was selected from my portfolio because it represents a transition from a standard academic assignment to an \textbf{enterprise-grade SaaS solution}. It demonstrates high-level proficiency in full-stack engineering, security, and DevOps.

\subsection{1. Core Problem Solved}
Traditional \textit{Umuganda} (community work) coordination relied on manual registration and fragmented communication. The system solves three critical architectural problems:
\begin{enumerate}[noitemsep]
    \item \textbf{Administrative Synchronization}: Centralizing localized village data into a national hierarchical oversight portal.
    \item \textbf{Geographic Intelligence}: Real-time spatial mapping across all administrative levels (District, Sector, Cell, Village).
    \item \textbf{Data Integrity}: Replacing paper-based records with automated attendance tracking and a secure system audit log.
\end{enumerate}

\subsection{2. Reporting \& Analytics Improvements}
The modernization phase introduced a professional reporting engine that significantly improved administrative oversight:
\begin{itemize}[noitemsep]
    \item \textbf{Branded PDF Automation}: Automated generation of board-ready activity summaries.
    \item \textbf{Data Export Ecosystem}: High-performance registry exports in CSV for offline longitudinal analysis.
    \item \textbf{Real-time Trends}: Visual representation of citizen engagement and attendance rates via Recharts.
\end{itemize}

\section{Technical Depth \& Architecture}
The project emphasizes technical rigor over simple functionality, addressing "Software Engineering" principles rather than basic "Web Development."

\begin{description}
    \item[Full-Stack Separation:] A clear decoupled architecture using Spring Boot (Backend) and React 19 (Frontend), communicating via documented REST APIs.
    \item[Security First:] Implementation of \textbf{Stateless JWT (JSON Web Token)} and \textbf{Role-Based Access Control (RBAC)} (Admin, Chef, Social, Villager).
    \item[Modern Tooling:] Leveraging \textbf{Vite} for optimized builds, \textbf{Tailwind CSS 4} for modern styling, and \textbf{Redux Toolkit} for predictable state management.
    \item[Stability \& DevOps:] Fully \textbf{Dockerized} environment (Postgres, API, Nginx) ensuring identical behavior across development and production. The live environment is powered directly by images hosted on \textbf{DockerHub}.
\end{description}

\section{Modernization Strategy (The "Problem Solver" Proof)}
To demonstrate independent problem-solving, I upgraded the baseline project to modern standards:
\begin{itemize}[noitemsep]
    \item \textbf{React 19 Transition}: Successfully ported the app to React 19 and Vite, achieving a significant reduction in bundle size and 40\% faster HMR.
    \item \textbf{PWA Implementation}: Converted the app into a Progressive Web App to support offline capabilities for field leaders.
    \item \textbf{Performance Optimization}: Implemented Skeleton UI architectures to manage asynchronous data loading smoothly.
\end{itemize}

\newpage

\section{Loom Presentation Script (4 Minutes)}
The goal of this presentation is to highlight engineering decisions rather than just clicking features.

\subsection{Phase 1: Introduction (0:00 - 0:30)}
\begin{itemize}
    \item \textbf{Visual:} Landing Page.
    \item \textbf{Narrative:} Introducing the Umuganda System and its role in modernizing community coordination in Rwanda.
\end{itemize}

\subsection{Phase 2: The Problem \& UI (0:30 - 1:15)}
\begin{itemize}
    \item \textbf{Visual:} Admin Dashboard with Mapping UI.
    \item \textbf{Narrative:} Discussing geographic synchronization and showing the integration of \textbf{Leaflet.js} and \textbf{Tailwind CSS 4}.
\end{itemize}

\subsection{Phase 3: Architecture Deep Dive (1:15 - 2:15)}
\begin{itemize}
    \item \textbf{Visual:} \texttt{SecurityConfig.java} snippet or Swagger UI.
    \item \textbf{Narrative:} Explaining the stateless JWT security, RBAC hierarchy, and the \textbf{Modular Monolith} backend strategy.
\end{itemize}

\subsection{Phase 4: DevOps \& Challenges (2:15 - 3:15)}
\begin{itemize}
    \item \textbf{Visual:} \texttt{Dockerfile} or System Audit Log.
    \item \textbf{Narrative:} Demonstrating Docker containerization and the System Audit Trail for accountability.
\end{itemize}

\subsection{Phase 5: Conclusion (3:15 - 4:00)}
\begin{itemize}
    \item \textbf{Visual:} Contact Info/Login.
    \item \textbf{Narrative:} Summary of technical impact and personal contributions to this "Enterprise-Grade" project.
\end{itemize}

\section{Final Email Draft}
\textit{Template to be used for the direct reply to Ferdinand.}

\begin{quote}
Dear Ferdinand,

Thank you for the opportunity to share my technical work. I am pleased to submit my most recent project: the \textbf{Umuganda Management System.} 

Originally started as an academic assignment, I independently modernized the system to enterprise SaaS standards. It is a full-stack platform built with \textbf{Spring Boot 3} and \textbf{React 19} designed to digitalize and secure community coordination across Rwanda.

\textbf{Key Technical Features:}
\begin{itemize}
    \item \textbf{Architecture}: Decoupled REST API (Java/Maven) with a Vite-optimized React frontend.
    \item \textbf{Security}: Stateless JWT authentication with extensive Role-Based Access Control (RBAC).
    \item \textbf{Geographic Intelligence}: Interactive mapping using Leaflet.js for administrative oversight.
    \item \textbf{DevOps}: Fully containerized using Docker for seamless deployment.
\end{itemize}

\textbf{Submission Links:}
\begin{itemize}
    \item \textbf{Code Repository}: \url{https://github.com/andremugabo/umuganda_system}
    \item \textbf{Live Demo (Frontend)}: \url{https://umuganda-frontend.onrender.com} (Docker image from DockerHub)
    \item \textbf{Interactive API Docs}: \url{https://umuganda-backend-k32m.onrender.com/swagger-ui/index.html} (Docker image from DockerHub)
    \item \textbf{Demo Admin Credentials}: \texttt{mugaboandre@umuganda.rw} / \texttt{123456}
\end{itemize}

I have also attached a brief video walkthrough [Loom Link] where I dive deeper into the engineering decisions and challenges I solved during the modernization phase.

Best regards, \\
Andre Mugabo
\end{quote}

\end{document}
